Задача трех тел, одна из классических задач небесной механики, описывается системой уравнений
$$\ddot q_i = \sum_{j \neq i} G m_j \frac{q_i-q_j}{|q_i-q_j|^3}, \qquad i = 0,1,2,$$
где $G$ -- гравитационная постоянная, $m_i$ -- массы тел.
Предполая, что $m_2 \to 0$ получаем \textit{ограниченную} задачу трех тел:
$$\ddot q_0 = m_1 \frac{q_1-q_0}{|q_1-q_0|^3}, \qquad \ddot q_1 = m_0 \frac{q_0-q_1}{|q_0-q_1|^3}, \qquad \ddot q_2 = m_0 \frac{q_0-q_2}{|q_0-q_2|^3} + m_1 \frac{q_1-q_2}{|q_1-q_2|^3}.$$
Первые два уравнения описывают Кеплеровскую систему из двух массивных тел (например \textit{Солнца} и \textit{Юпитера}) с массами $m_0$ и $m_1$, вращающихся вокруг общего центра масс по эллиптичкской орбите. Третье уравнение описывает движение \textit{безмассового} тела (будем называть его \textit{астероидом}) в гравитационном поле массивных тел, не оказывая на них влияния.

Если движение безмассового тела происходит в плоскости движения массивных тел, задача становится плосокй и как правило называется плоской эллиптической ограниченной задачей трёх тел (далее ПЭОЗТ). Такая аппроксимация имеет множество практических применений. Одно из них связано с описанием движения в главном поясе астероиодов -- области в солнечной системе между орбитами Марса и Юпитера, содержащей миллионы объектов размером от пылевых частиц до астероидов. Пренебрегая всеми прочими телами (кроме Солнца и Юпитера) ввиду малости их масс и большого удаления, динамику объектов в поясе  астероидов можно описывать ПЭОЗТ.

Обозначим за $\nu = m_1/(m_0+m_1)$ отношение масс массивных тел, где $m_0$ -- масса Солнца, а $m_1$ -- масса Юпитера. Так как $\nu \approx 1/1047$, влияние Юпитера на движение астероидов мало, поэтому орбиты астероидов близки к эллиптическим. Однако в 1866 году Д. Кирквудом было замечено, что распредление больших полуосей этих орбит содержит ямы (позже названные щели Кирквуда). Кроме того было обнаружено, что этим ямам соответствуют орбитальные резонансы, то есть целочисленные отношения периода Юпитера и периодов астероидов.

В 1980-х годах вышла серия работ Д. Висдома, посвященная изучению механизма формирования щелей Кирквуда. Используя ПЭОЗТ он исследовал динамику астероидов в присутствии орбитального резонанса, в частности в случае резонанса 3:1, cвязанного с самой заметной щелью Кирквуда в поясе астероидов.
В работе \cite{wis1} он ввел модельную гамильтонову систему, которая дальше будет называться системой Висдома. 
Следует отметить, что данная система представляет интерес не только из-за связи с поясом астероидов, а также потому, что является допустимой аппроксимацией для движения астероида в поле двух произвольным массивных тел отношение масс которых мало.

Одной из ключевых особенностей системы Висдома является наличие двух различных временных масштабов в той динамике, которую она проявляет. Это позволяет выделить пару ''быстрых'' и пару ''медленных'' переменных. Причем при фиксированных ''медленных'' переменных, динамика ''быстрых'' описывается уравнением типа маятника.

Система Висдома, будучи гамильтоновой системой с двумя степенями свободы, исследовалась в ряде работ. В частности, в !!![9] был разработан полуаналитический метод изучения долговременной динамики в окрестности резонанса 3:1. С помощью этого подхода была обнаружена обширная хаотическая зона вблизи указанного резонанса, для которого характерны траектории со значительным и быстрым изменением эксцентриситетов !!![10, 11]. Численное интегрирование показало, что такие хаотические траектории могут достигать значений эксцентриситетов, достаточных для пересечения орбит Марса и Земли. В работе !!![11] подчеркивается, что одной из возможных причин возникновения подобных хаотических режимов является пересечение сепаратрисы, т.е. выход из области применимости адиабатического приближения. Полный анализ динамики ''медленной'' переменной представлен в !!![6]. Кроме того, численно подтверждено существование островков устойчивости в зоне адиабатического хаоса.

Целью данной работы является изложение некоторых результатов, касающихся системы Висдома, которая рассматривается как быстро-медленная гамильтонова система с двумя быстрыми и двумя медленными переменными. Основное внимание уделяется геометрии ее медленного многообразия, а также динамике системы Висдома в окрестности этого многообразия. Особое внимание уделяется анализу пересечений инвариантных многообразий, связанных с положениями равновесия системы.

\section{Определение системы Висдома}

В своей диссертации !!!![9] Дж. Висдом рассматривает следующее приближение для резонансного гамильтониана в окрестности резонанса 3:1 в ПЭОЗТ:
\begin{equation}
    H_R = - \frac{(1-\nu)^2}{2L^2} - \nu R_{sec}(\rho, \omega) - \nu R_{res}(L,l,\rho,\omega).
    \label{hamres}
\end{equation}
Здесь $R_{sec}$ и  $R_{res}$ -- секулярная и резонансная части возмущающей функции соответственно:
$$R_{sec}(\rho, \omega) = - 2 \rho F - e_J G \sqrt{2 \rho} \cos \omega,$$
$$R_{res}(L,l,\rho,\omega) = 2 \rho C \cos(l-\omega-3 l_J)+e_J^2E \cos(l - 3 l_J).$$

В координатах Делоне-Пуанкаре импульсы $L$ и $l$ могут быть выражены через большую полуось орбиты астероида $a$ и ее эксцентриситет $e$:
$$L = \sqrt{(1-\nu)a}, \qquad \sqrt{(1-\nu)a} (1-(1-e^2)^{1/2}).$$

Сопряженными к ним координатами являются средняя долгота астероида $l$ и взятая с обратным знаком долгота его перигелия $\omega$. Выбор единиц измерения осуществляется таким образом, чтобы большая полуось Юпитера $a_J$ и его средняя долгота $l_J$ удовлетворяли условиям 
$$a_J=1, l_J= t .$$
Как следствие период Юпитера становится равным $2\pi$.
Коэффициенты $C,D,E,F,G$ зависят от резонансной большой полуоси астероида $a_{res} = ((1-\nu)/9)^{1/3}$. В частности для $\nu = 1/1047.355$
$$C \approx 0.8631579, \quad D \approx -2.656407, \quad E \approx 0.3629536, \quad F \approx -0.2050694, \quad G \approx 0.1987054.$$

Следуя работе !!!![9], выполним каноническую (относительно стандартной симплектической структуры $\Omega_{st}=dL \wedge dl+d \rho \wedge d \omega$) замену переменных с помощью зависящей от времени производящей функции:
$$S=(l-3t)(\Lambda+L_{res})+\frac{x^2}2 \tan \omega$$
и разложим (\ref{hamres}) с точностью до квадратичных членов по $\Lambda$ в окрестности резонансного значения $L_{res}$, определяемого условием
$$\frac{\partial H}{\partial \Lambda}\bigg|_{\Lambda=0} = 0.$$

Это приводит гамильтониан к следующему виду !!!![9]:
$$H = \alpha \frac{\Lambda^2}{2} + \nu \big( H_0(e_Jx, e_Jy) + H_1(\lambda, e_Jx, e_J y) \big),$$
где
$$H_0(x,y) = Gx +F(x^2+y^2), \qquad H_1(\lambda,x,y) = U(x,y) \cos \lambda + V(x,y) \sin \lambda,$$
$$U(x,y) = C(x^2-y^2) + Dx + E, \qquad V(x,y) = 2Cxy + Dy,$$
$$\alpha = - \frac{3 (1-\nu)^2}{L_{res}^4} < 0, \qquad L_{res} = \left( \frac{(1-\nu)^2}{3} \right)^{1/3}.$$


Наконец, введем новый малый параметр $ \varepsilon = \sqrt{\nu} $ и выполним неканоническую перенормировку переменных и параметра $ \alpha $:

$$t \rightarrow \varepsilon t,\quad \Lambda \rightarrow (\varepsilon e_J^2)^{-1} \Lambda,\quad x \rightarrow e_J^{-1} x,\quad y \rightarrow e_J^{-1} y,\quad \alpha \rightarrow e_J^2 \alpha$$

В новых переменных (которые мы обозначаем прежними буквами) уравнения движения принимают вид гамильтоновой системы с функцией Гамильтона

\begin{equation}
    H_w = \alpha \frac{\Lambda^2}{2} + H_0(x, y) + H_1(\lambda, x, y),
    \label{hw}
\end{equation}
относительно зависящей от $ \varepsilon $ симплектической формы

\begin{equation}
\Omega_w = d\Lambda \wedge d\lambda + \varepsilon^{-1} dx \wedge dy.
\label{ow}
\end{equation}

Данная гамильтонова система является объектом нашего исследования; далее мы будем называть ее (следуя работе !!![6]) системой Висдома. В силу $ 2\pi $-периодичности $ H_w $ по переменной $ \lambda $, фазовое пространство системы Висдома представляет собой $ S^1 \times \mathbb{R}^3 $.

Примечательно, что в работе [6] было замечено, что гамильтониан $ H_w $ не зависит от эксцентриситета Юпитера $ e_J $. Хотя последний и определяет величины вековых изменений элементов орбиты астероида, в главном приближении (при $ e_J \rightarrow 0+ $) он не влияет на топологические свойства фазового потока для ПЭОЗТ, усредненной в окрестности резонанса.

В настоящей статье мы исследуем непосредственно систему Висдома, рассматривая ее вне связи с динамикой в поясе астероидов. Будем считать, что $ \alpha, B, C, D, E, F $ некоторые вещественные константы, удовлетворяющие следующему

\begin{hyp}
Будем считать, что константы $C,D,E,F,G$ удовлетворяют следующим условиям:
\begin{subequations}
\begin{gather}
D^2 - 4CE > 0, \label{hypb} \\
C > |F|, \label{hypa} \\
(DF - GC)^2 < (D^2 - 4CE) (|F| - C)^2. \label{hypc}
\end{gather}
\end{subequations}
\end{hyp}


Все дополнительные предположения относительно этих констант будут оговариваться особо по ходу изложения.